\documentclass[aspectratio=169,11pt]{beamer}

% --- Theme ---
\usetheme{metropolis}
\usepackage{appendixnumberbeamer}

% --- Packages ---
\usepackage{fontspec}
\usepackage{minted}
\setminted{fontsize=\small,breaklines=true}
\usepackage{booktabs}
\usepackage{tikz}
\usepackage{fontawesome5}

% --- Colors ---
\definecolor{healthgreen}{RGB}{34,139,34}
\definecolor{alertred}{RGB}{220,53,69}
\setbeamercolor{alerted text}{fg=alertred}

% --- Metadata ---
\title{Module 6: Hypothesis testing for health data}
\subtitle{Data analysis with Python for health specialists}
\author{Ya\'{e} Ulrich Gaba}
\institute{AIRINA Labs}
\date{2026}

\begin{document}

% ============================================================
\maketitle

% ============================================================
\begin{frame}{The logic of hypothesis testing}
Every clinical question $\to$ two competing statements:
\begin{itemize}
  \item $H_0$ (null): There is \alert{no} effect / no difference / no association.
  \item $H_1$ (alternative): There \alert{is} an effect / a difference / an association.
\end{itemize}

\vspace{6pt}
\textbf{The 5 steps:}
\begin{enumerate}
  \item State hypotheses ($H_0$ and $H_1$)
  \item Choose significance level $\alpha$ (usually 0.05)
  \item Collect data, compute test statistic
  \item Compute p-value
  \item Decide: if $p < \alpha$, reject $H_0$
\end{enumerate}

\begin{exampleblock}{Clinical framing}
A statistician writes $H_0: \mu_1 = \mu_2$. A clinician writes ``no difference in SBP between statin and placebo groups.'' Both say the same thing.
\end{exampleblock}
\end{frame}

% ============================================================
\begin{frame}{p-values: what they mean and what they don't}
\textbf{A p-value IS:} the probability of seeing data this extreme if $H_0$ is true.

\vspace{6pt}
\textbf{A p-value is NOT:}
\begin{itemize}
  \item The probability that $H_0$ is true
  \item The probability the result is ``due to chance''
  \item A measure of effect size
  \item Meaningfully different at 0.049 vs.\ 0.051
\end{itemize}

\vspace{6pt}
\begin{alertblock}{The big mistake}
A study with 100\,000 patients can find $p < 0.001$ for a BP difference of 0.5~mmHg --- \alert{statistically significant but clinically meaningless}. Always report \textbf{effect size} alongside p-values.
\end{alertblock}
\end{frame}

% ============================================================
\begin{frame}[fragile]{One-sample t-test}
\textbf{Q:} Is mean cholesterol in our clinic different from the national average of 200~mg/dL?

\begin{minted}{python}
from scipy import stats
import numpy as np

np.random.seed(42)
cholesterol = np.random.normal(loc=212, scale=35, size=80)

t_stat, p_value = stats.ttest_1samp(cholesterol, popmean=200)
print(f"Sample mean: {cholesterol.mean():.1f} mg/dL")
print(f"t = {t_stat:.3f}, p = {p_value:.4f}")
\end{minted}

Compares a \textbf{sample mean} to a \textbf{known reference value} (guideline threshold, national average).
\end{frame}

% ============================================================
\begin{frame}[fragile]{Two-sample t-test}
\textbf{Q:} Is systolic BP different between treatment and control groups?

\begin{minted}{python}
np.random.seed(42)
treatment = np.random.normal(loc=128, scale=15, size=60)
control = np.random.normal(loc=138, scale=16, size=55)

# Welch's t-test (safer default: no equal variance assumption)
t_stat, p_value = stats.ttest_ind(treatment, control,
                                   equal_var=False)
print(f"Difference: {control.mean()-treatment.mean():.1f} mmHg")
print(f"p = {p_value:.4f}")
\end{minted}

\begin{exampleblock}{Check assumptions}
Normality: \texttt{stats.shapiro()}. Equal variances: \texttt{stats.levene()}. If variances differ, use Welch's t-test (\texttt{equal\_var=False}).
\end{exampleblock}
\end{frame}

% ============================================================
\begin{frame}[fragile]{Paired t-test}
\textbf{Q:} Does a 12-week exercise program reduce blood pressure?

\begin{minted}{python}
np.random.seed(42)
bp_before = np.random.normal(loc=142, scale=12, size=45)
bp_after = bp_before - np.random.normal(loc=8, scale=6, size=45)

t_stat, p_value = stats.ttest_rel(bp_before, bp_after)
diffs = bp_before - bp_after
print(f"Mean reduction: {diffs.mean():.1f} mmHg")
print(f"p = {p_value:.6f}")
\end{minted}

\begin{alertblock}{Common mistake}
Using an \textbf{independent} t-test on before/after data ignores the pairing and loses statistical power. Same patients measured twice $\to$ always use a \textbf{paired} t-test.
\end{alertblock}
\end{frame}

% ============================================================
\begin{frame}[fragile]{Chi-square test of independence}
\textbf{Q:} Is diabetes prevalence associated with obesity category?

\begin{minted}{python}
# Build contingency table
contingency = pd.crosstab(df["bmi_category"],
                          df["diabetes_label"])

# Chi-square test
chi2, p_value, dof, expected = stats.chi2_contingency(
    contingency)
print(f"Chi2 = {chi2:.2f}, dof = {dof}, p = {p_value:.4f}")
\end{minted}

Works with \textbf{categorical} variables. Requires expected cell counts $\geq 5$. For small cells, use Fisher's exact test: \texttt{stats.fisher\_exact(table)}.
\end{frame}

% ============================================================
\begin{frame}[fragile]{Mann-Whitney U test}
When data is \alert{not normally distributed} (length of stay, costs, small samples):

\begin{minted}{python}
np.random.seed(42)
los_surgical = np.random.exponential(scale=7, size=40)
los_medical = np.random.exponential(scale=4.5, size=45)

u_stat, p_value = stats.mannwhitneyu(
    los_surgical, los_medical, alternative="two-sided")
print(f"Median surgical: {np.median(los_surgical):.1f} days")
print(f"Median medical:  {np.median(los_medical):.1f} days")
print(f"p = {p_value:.4f}")
\end{minted}

\textbf{Rule of thumb:} Normal $\to$ t-test. Non-normal $\to$ Mann-Whitney. Categorical $\to$ chi-square. Paired non-normal $\to$ Wilcoxon signed-rank.
\end{frame}

% ============================================================
\begin{frame}[fragile]{Multiple testing correction}
Testing 20 hypotheses at $\alpha = 0.05$ $\to$ expect 1 false positive even if nothing is going on.

\begin{minted}{python}
from statsmodels.stats.multitest import multipletests

# Bonferroni: divide alpha by number of tests
rejected_bonf, _, _, _ = multipletests(
    p_values, alpha=0.05, method="bonferroni")

# FDR (Benjamini-Hochberg): less conservative
rejected_fdr, _, _, _ = multipletests(
    p_values, alpha=0.05, method="fdr_bh")
\end{minted}

\textbf{Bonferroni:} strict, minimizes false positives (confirmatory trials).\\
\textbf{FDR:} less strict, controls false discovery rate (exploratory/genomics).
\end{frame}

% ============================================================
\begin{frame}[fragile]{Effect size: Cohen's d}
$$d = \frac{\bar{x}_1 - \bar{x}_2}{s_{\text{pooled}}}$$

\begin{itemize}
  \item $|d| < 0.2$: negligible
  \item $0.2 \leq |d| < 0.5$: small
  \item $0.5 \leq |d| < 0.8$: medium
  \item $|d| \geq 0.8$: large
\end{itemize}

\begin{minted}{python}
def cohens_d(g1, g2):
    n1, n2 = len(g1), len(g2)
    pooled = np.sqrt(((n1-1)*g1.var(ddof=1) +
                      (n2-1)*g2.var(ddof=1)) / (n1+n2-2))
    return (g1.mean() - g2.mean()) / pooled
\end{minted}

\alert{Always report effect size alongside p-values.} Statistical significance $\neq$ clinical significance.
\end{frame}

% ============================================================
\begin{frame}{What you can do after this module}
\begin{enumerate}
  \item Formulate $H_0$/$H_1$ for clinical research questions
  \item Apply the right test: t-test (1-sample, 2-sample, paired), chi-square, Mann-Whitney
  \item Interpret p-values correctly and avoid common misconceptions
  \item Correct for multiple comparisons (Bonferroni, FDR)
  \item Compute Cohen's $d$ to assess clinical meaningfulness
\end{enumerate}

\vspace{12pt}
\centering
\Large \alert{Next: Module 7 --- Regression for health outcomes}
\end{frame}

% ============================================================
\begin{frame}{Exercises (Hour 3)}
\begin{enumerate}
  \item \textbf{One-sample t-test:} Framingham cholesterol vs.\ 200 mg/dL
  \item \textbf{Two-sample t-test:} Systolic BP in CHD vs.\ non-CHD groups with Cohen's $d$
  \item \textbf{Chi-square:} Smoking vs.\ 10-year CHD risk
  \item \textbf{Mini-project:} Run all tests on the Framingham dataset with multiple testing correction; identify the variable with the largest effect size
\end{enumerate}
\end{frame}

\end{document}
